\section{Eventually different families over stationary supports}

\import{silver/power-closure.tex}
\import{stationary/fodor.tex}
\import{cardinal/families.tex}

\subsection{Agreement and stationary supports}

\begin{definition}\label{silver_eventual_equality_set}
    $\silverEventualEqualitySet{f}{g}{\lambda}
        =\{i\in\lambda\mid f(i)=g(i)\}$.
\end{definition}

\begin{definition}\label{eventually_different}
    $\eventuallyDifferent{f}{g}{\lambda}$ iff
        $\ordinalBounded{\silverEventualEqualitySet{f}{g}{\lambda}}
            {\lambda}$.
\end{definition}

\begin{definition}\label{eventually_different_family}
    $\eventuallyDifferentFamily{F}{\lambda}$ iff
        for all $f,g\in F$ such that $f\neq g$ we have
            $\eventuallyDifferent{f}{g}{\lambda}$.
\end{definition}

\begin{proposition}\label{eventually_different_family_subset}
    Suppose $\eventuallyDifferentFamily{F}{\lambda}$.
    Let $G$ be a subset of $F$.
    Then $\eventuallyDifferentFamily{G}{\lambda}$.
\end{proposition}
\begin{proof}
    Follows by \obv{eventually_different_family,subseteq}.
\end{proof}

\begin{proposition}\label{silver_eventual_equality_set_subseteq}
    $\silverEventualEqualitySet{f}{g}{\lambda}\subseteq\lambda$.
\end{proposition}
\begin{proof}
    Follows by \obv{silver_eventual_equality_set,subseteq}.
\end{proof}

\begin{proposition}\label{eventually_different_agreement_set_nonstationary}
    Let $\lambda$ be a limit ordinal.
    Suppose $S\subseteq\lambda$.
    Suppose $\eventuallyDifferent{f}{g}{\lambda}$.
    Suppose for every $i\in S$ we have $f(i)=g(i)$.
    Then $\nonstationaryIn{S}{\lambda}$.
\end{proposition}
\begin{proof}
    Take $\beta\in\lambda$ such that
        $\silverEventualEqualitySet{f}{g}{\lambda}\subseteq\beta$
        by \auto{eventually_different,ordinal_bounded}.
    $S\subseteq\silverEventualEqualitySet{f}{g}{\lambda}$
        by \obv{subseteq,silver_eventual_equality_set}.
    Hence $S\subseteq\beta$ by \obv{subseteq_transitive}.
    Hence $\ordinalBounded{S}{\lambda}$
        by \obv{ordinal_bounded}.
    Follows by \obv{bounded_is_nonstationary}.
\end{proof}

We restrict the support to limit indices of the normal sequence, where
continuity gives regressive ranks.

\begin{definition}\label{silver_limit_support}
    $\silverLimitSupport{T}{\lambda}
        =T\inter\limitOrdinalsBelow{\lambda}$.
\end{definition}

\begin{theorem}\label{silver_limit_support_is_stationary}
    Let $\lambda$ be a regular cardinal.
    Suppose $\naturals\in\lambda$.
    Suppose $T$ is stationary in $\lambda$.
    Then $\silverLimitSupport{T}{\lambda}$ is stationary in $\lambda$.
\end{theorem}
\begin{proof}
    $\lambda$ is an infinite cardinal and
        $\cofinality{\lambda}=\lambda$ by \obv{regular_cardinal}.
    Hence $\naturals\in\cofinality{\lambda}$.
    $\clubIn{\limitOrdinalsBelow{\lambda}}{\lambda}$
        by \obv{limit_ordinals_below_is_club}.
    $T\inter\limitOrdinalsBelow{\lambda}$ is stationary in $\lambda$
        by \obv{stationary_intersection_club}.
    Follows by \obv{silver_limit_support}.
\end{proof}

\subsection{Pressing down the coordinate ranks}

Choose coordinate injections $E(i)$ from $A(i)$ into $c(i)$ on the support.
At a limit index $i$, the ordinal $(E(i))(f(i))$ belongs to some $c(j)$
with $j<i$. The rank of $f$ at $i$ is the least such $j$.

\begin{definition}\label{silver_rank_candidates}
    $\silverRankCandidates{E}{c}{f}{i}
        =\{j\in i\mid (E(i))(f(i))\in c(j)\}$.
\end{definition}

\begin{definition}\label{silver_rank_map}
    $\silverRankMap{E}{c}{f}{T}
        =\{(i,\inters{\silverRankCandidates{E}{c}{f}{i}})
            \mid i\in T\}$.
\end{definition}

\begin{proposition}\label{silver_rank_candidates_inhabited}
    Suppose $\normalMap{c}{\lambda}{\kappa}$.
    Suppose $T\subseteq\limitOrdinalsBelow{\lambda}$.
    Suppose $A$ is a function on $\lambda$.
    Suppose $E$ is a function on $T$.
    Suppose for every $j\in T$ we have
        $E(j)\in\inj{A(j)}{c(j)}$.
    Let $f\in\dependentProduct{\lambda}{A}$.
    Let $i\in T$.
    Then $\silverRankCandidates{E}{c}{f}{i}$ is inhabited.
\end{proposition}
\begin{proof}
    $i\in\limitOrdinalsBelow{\lambda}$ by \obv{subseteq}.
    Hence $i\in\lambda$ by \obv{limit_ordinals_below}.
    Hence $i$ is a limit ordinal by \obv{limit_ordinals_below}.
    $\lambda$ is an ordinal by \auto{normal_map_type}.
    $c\in\funs{\lambda}{\kappa}$
        by \obv{normal_map_type}.
    Hence $c$ is a function and $\dom{c}=\lambda$
        by \obv{funs_is_function,funs}.
    $f(i)\in A(i)$ by \obv{dependent_product_coordinate}.
    $E(i)\in\inj{A(i)}{c(i)}$.
    Hence $(E(i))(f(i))\in c(i)$ by \auto{inj,funs_type_apply}.
    $c(i)=\unions{\ran{\restrl{c}{i}}}$
        by \obv{normal_map_continuous_elim}.
    Let $r=\restrl{c}{i}$.
    Hence $c(i)=\unions{\ran{r}}$.
    Hence $(E(i))(f(i))\in\unions{\ran{r}}$.
    Take $z\in\ran{r}$ such that $(E(i))(f(i))\in z$
        by \obv{unions_iff}.
    $r$ is a function by \obv{restrl_of_function_is_function}.
    $i\subseteq\lambda$ by \obv{ordinal_elem_implies_subseteq}.
    Hence $\dom{r}=i$
        by \auto{dom_of_restrl_eq_inter,inter_absorb_supseteq_left}.
    Take $j\in\dom{r}$ such that $r(j)=z$
        by \obv{fun_ran_iff}.
    Hence $j\in i$.
    $r(j)=c(j)$ by \auto{restrl_of_function_apply}.
    Hence $(E(i))(f(i))\in c(j)$.
    Hence $j\in\silverRankCandidates{E}{c}{f}{i}$
        by \obv{silver_rank_candidates}.
\end{proof}

\begin{proposition}\label{silver_rank_value_properties}
    Suppose $\normalMap{c}{\lambda}{\kappa}$.
    Suppose $T\subseteq\limitOrdinalsBelow{\lambda}$.
    Suppose $A$ is a function on $\lambda$.
    Suppose $E$ is a function on $T$.
    Suppose for every $j\in T$ we have
        $E(j)\in\inj{A(j)}{c(j)}$.
    Let $f\in\dependentProduct{\lambda}{A}$.
    Let $i\in T$.
    Let $R=\silverRankCandidates{E}{c}{f}{i}$.
    Then $\inters{R}\in i$ and $(E(i))(f(i))\in c(\inters{R})$.
\end{proposition}
\begin{proof}
    $R$ is inhabited by \auto{silver_rank_candidates_inhabited}.
    $i\in\limitOrdinalsBelow{\lambda}$ by \obv{subseteq}.
    Hence $i$ is a limit ordinal by \obv{limit_ordinals_below}.
    Hence $i$ is an ordinal by \obv{limit_ordinal}.
    For every $j\in R$ we have $j\in i$
        by \auto{silver_rank_candidates}.
    Hence for every $j\in R$ we have $j$ is an ordinal
        by \obv{elem_of_ordinal_is_ordinal}.
    $\inters{R}\in R$ by \obv{inters_of_ordinals_elem}.
    Hence $\inters{R}\in i$ and $(E(i))(f(i))\in c(\inters{R})$
        by \auto{silver_rank_candidates}.
\end{proof}

\begin{theorem}[The rank map is regressive]\label{silver_rank_map_type_and_regressive}
    Suppose $\normalMap{c}{\lambda}{\kappa}$.
    Suppose $T\subseteq\limitOrdinalsBelow{\lambda}$.
    Suppose $A$ is a function on $\lambda$.
    Suppose $E$ is a function on $T$.
    Suppose for every $j\in T$ we have
        $E(j)\in\inj{A(j)}{c(j)}$.
    Let $f\in\dependentProduct{\lambda}{A}$.
    Let $r=\silverRankMap{E}{c}{f}{T}$.
    Then $r\in\funs{T}{\lambda}$ and $r$ is regressive on $T$.
\end{theorem}
\begin{proof}
    $\lambda$ is an ordinal by \auto{normal_map_type}.
    $r$ is a relation by \auto{relation,silver_rank_map}.
    $r$ is right-unique
        by \auto{rightunique,silver_rank_map,pair_eq_iff}.
    $\dom{r}=T$
        by \auto{setext,dom_iff,silver_rank_map,pair_eq_iff}.
    Show that for every $i\in T$ we have $r(i)\in i$ and
        $r(i)\in\lambda$.
    \begin{subproof}
        Fix $i\in T$.
        Let $R=\silverRankCandidates{E}{c}{f}{i}$.
        $\inters{R}\in i$ by \auto{silver_rank_value_properties}.
        $(i,\inters{R})\in r$ by \obv{silver_rank_map}.
        Hence $r(i)=\inters{R}$ by \obv{function_apply_intro}.
        $i\in\limitOrdinalsBelow{\lambda}$ by \obv{subseteq}.
        Hence $i\in\lambda$ by \obv{limit_ordinals_below}.
        Hence $r(i)\in\lambda$ by \obv{ordinal_transitivity}.
    \end{subproof}
    Hence $r$ is a function to $\lambda$.
    Hence $r\in\funs{T}{\lambda}$ by \obv{funs_intro}.
    For every $i\in T\setminus\{\emptyset\}$ we have $r(i)\in i$
        by \obv{setminus_elim_left}.
    Hence $r$ is regressive on $T$ by \obv{regressive_on}.
\end{proof}

\begin{proposition}\label{silver_rank_map_codes_value}
    Suppose $\normalMap{c}{\lambda}{\kappa}$.
    Suppose $T\subseteq\limitOrdinalsBelow{\lambda}$.
    Suppose $A$ is a function on $\lambda$.
    Suppose $E$ is a function on $T$.
    Suppose for every $j\in T$ we have
        $E(j)\in\inj{A(j)}{c(j)}$.
    Let $f\in\dependentProduct{\lambda}{A}$.
    Let $i\in T$.
    Let $r=\silverRankMap{E}{c}{f}{T}$.
    Then $(E(i))(f(i))\in c(r(i))$.
\end{proposition}
\begin{proof}
    $r$ is a function by
        \auto{silver_rank_map_type_and_regressive,funs_is_function}.
    Let $R=\silverRankCandidates{E}{c}{f}{i}$.
    $(i,\inters{R})\in r$ by \obv{silver_rank_map}.
    Hence $r(i)=\inters{R}$ by \obv{function_apply_intro}.
    $(E(i))(f(i))\in c(\inters{R})$
        by \auto{silver_rank_value_properties}.
\end{proof}

Choose the least rank whose fibre is stationary; that fibre will be the
support of the code.

\begin{definition}\label{silver_stationary_ranks}
    $\silverStationaryRanks{E}{c}{f}{T}{\lambda}
        =\{j\in\lambda\mid
            \text{$\fibre{\silverRankMap{E}{c}{f}{T}}{j}$ is stationary in
                $\lambda$}\}$.
\end{definition}

\begin{definition}\label{silver_rank_index}
    $\silverRankIndex{E}{c}{f}{T}{\lambda}
        =\inters{\silverStationaryRanks{E}{c}{f}{T}{\lambda}}$.
\end{definition}

\begin{definition}\label{silver_support}
    $\silverSupport{E}{c}{f}{T}{\lambda}
        =\fibre{\silverRankMap{E}{c}{f}{T}}
            {\silverRankIndex{E}{c}{f}{T}{\lambda}}$.
\end{definition}

\begin{theorem}\label{silver_stationary_ranks_inhabited}
    Let $\lambda$ be a regular cardinal.
    Suppose $\naturals\in\lambda$.
    Suppose $\normalMap{c}{\lambda}{\kappa}$.
    Suppose $T\subseteq\limitOrdinalsBelow{\lambda}$.
    Suppose $T$ is stationary in $\lambda$.
    Suppose $A$ is a function on $\lambda$.
    Suppose $E$ is a function on $T$.
    Suppose for every $j\in T$ we have
        $E(j)\in\inj{A(j)}{c(j)}$.
    Let $f\in\dependentProduct{\lambda}{A}$.
    Then $\silverStationaryRanks{E}{c}{f}{T}{\lambda}$ is inhabited.
\end{theorem}
\begin{proof}
    Let $r=\silverRankMap{E}{c}{f}{T}$.
    $r\in\funs{T}{\lambda}$ and $r$ is regressive on $T$
        by \obv{silver_rank_map_type_and_regressive}.
    Take $j\in\lambda$ such that
        $\fibre{r}{j}$ is stationary in $\lambda$ by \obv{fodor}.
    $j\in\silverStationaryRanks{E}{c}{f}{T}{\lambda}$
        by \obv{silver_stationary_ranks}.
\end{proof}

\begin{theorem}[A stationary support with bounded coordinate values]\label{silver_support_properties}
    Let $\lambda$ be a regular cardinal.
    Suppose $\naturals\in\lambda$.
    Suppose $\normalMap{c}{\lambda}{\kappa}$.
    Suppose $T\subseteq\limitOrdinalsBelow{\lambda}$.
    Suppose $T$ is stationary in $\lambda$.
    Suppose $A$ is a function on $\lambda$.
    Suppose $E$ is a function on $T$.
    Suppose for every $u\in T$ we have
        $E(u)\in\inj{A(u)}{c(u)}$.
    Let $f\in\dependentProduct{\lambda}{A}$.
    Let $m=\silverRankIndex{E}{c}{f}{T}{\lambda}$.
    Let $S=\silverSupport{E}{c}{f}{T}{\lambda}$.
    Then $m\in\lambda$ and $S\subseteq T$ and
        $S$ is stationary in $\lambda$ and
        for every $i\in S$ we have $(E(i))(f(i))\in c(m)$.
\end{theorem}
\begin{proof}
    $\lambda$ is a limit ordinal
        by \obv{infinite_cardinal_is_limit_ordinal}.
    Hence $\lambda$ is an ordinal by \obv{limit_ordinal}.
    Let $R=\silverStationaryRanks{E}{c}{f}{T}{\lambda}$.
    $R$ is inhabited by \auto{silver_stationary_ranks_inhabited}.
    For every $q\in R$ we have $q\in\lambda$
        by \auto{silver_stationary_ranks}.
    Hence for every $q\in R$ we have $q$ is an ordinal
        by \obv{elem_of_ordinal_is_ordinal}.
    $\inters{R}\in R$ by \obv{inters_of_ordinals_elem}.
    $m=\inters{R}$ by \auto{silver_rank_index}.
    Hence $m\in R$.
    Hence $m\in\lambda$ and
        $\fibre{\silverRankMap{E}{c}{f}{T}}{m}$ is stationary in $\lambda$
        by \obv{silver_stationary_ranks}.
    $S=\fibre{\silverRankMap{E}{c}{f}{T}}{m}$
        by \auto{silver_support}.
    Hence $S$ is stationary in $\lambda$.
    Let $r=\silverRankMap{E}{c}{f}{T}$.
    $r\in\funs{T}{\lambda}$
        by \auto{silver_rank_map_type_and_regressive}.
    Hence $r$ is a function and $\dom{r}=T$
        by \obv{funs_is_function,funs}.
    $S\subseteq\dom{r}$
        by \auto{subseteq,function_fibre_iff}.
    Hence $S\subseteq T$.
    Show that for every $i\in S$ we have $(E(i))(f(i))\in c(m)$.
    \begin{subproof}
        Fix $i\in S$.
        Hence $i\in\fibre{r}{m}$.
        Hence $i\in\dom{r}$ and $r(i)=m$
            by \obv{function_fibre_iff}.
        Hence $i\in T$.
        $(E(i))(f(i))\in c(r(i))$
            by \obv{silver_rank_map_codes_value}.
    \end{subproof}
\end{proof}

\subsection{Coding on the stationary support}

On the chosen support $S$, record the values $(E(i))(f(i))$ and extend by
zero off $S$. If the chosen rank is $m$, this gives a function from $\lambda$
to $c(m)$. The code records $m$, the support $S$, and this function.
Recording $S$ ensures that equal codes agree on the same stationary set.

\begin{definition}\label{silver_partial_code}
    $\silverPartialCode{E}{f}{S}
        =\{(i,(E(i))(f(i)))\mid i\in S\}$.
\end{definition}

\begin{definition}\label{silver_total_code}
    $\silverTotalCode{E}{f}{S}{\lambda}
        =\functionDomainExtension{\identity{S}}
            {\silverPartialCode{E}{f}{S}}{\lambda}{\emptyset}$.
\end{definition}

\begin{definition}\label{silver_code}
    $\silverCode{E}{c}{f}{T}{\lambda}
        =(\silverRankIndex{E}{c}{f}{T}{\lambda},
            (\silverSupport{E}{c}{f}{T}{\lambda},
             \silverTotalCode{E}{f}
                {\silverSupport{E}{c}{f}{T}{\lambda}}{\lambda}))$.
\end{definition}

\begin{definition}\label{silver_code_fibres}
    $\silverCodeFibres{c}{\lambda}
        =\{(j,\pow{\lambda}\times\funs{\lambda}{c(j)})
            \mid j\in\lambda\}$.
\end{definition}

\begin{definition}\label{silver_code_space}
    $\silverCodeSpace{c}{\lambda}
        =\indexedSum{\lambda}{\silverCodeFibres{c}{\lambda}}$.
\end{definition}

\begin{definition}\label{silver_code_map}
    $\silverCodeMap{E}{c}{F}{T}{\lambda}
        =\{(f,\silverCode{E}{c}{f}{T}{\lambda})\mid f\in F\}$.
\end{definition}

\begin{proposition}\label{silver_partial_code_type}
    Suppose for every $u\in S$ we have $(E(u))(f(u))\in B$.
    Let $p=\silverPartialCode{E}{f}{S}$.
    Then $p\in\funs{S}{B}$ and
        for every $v\in S$ we have $p(v)=(E(v))(f(v))$.
\end{proposition}
\begin{proof}
    $p$ is a relation by \auto{relation,silver_partial_code}.
    $p$ is right-unique
        by \auto{rightunique,silver_partial_code,pair_eq_iff}.
    $\dom{p}=S$
        by \auto{setext,dom_iff,silver_partial_code,pair_eq_iff}.
    For every $i\in S$ we have $(i,(E(i))(f(i)))\in p$
        by \obv{silver_partial_code}.
    Hence for every $i\in S$ we have $p(i)=(E(i))(f(i))$
        by \obv{function_apply_intro}.
    Hence for every $i\in S$ we have $p(i)\in B$ by \auto.
    Hence $p$ is a function to $B$.
    Hence $p\in\funs{S}{B}$ by \obv{funs_intro}.
\end{proof}

\begin{proposition}\label{silver_code_fibres_type}
    Let $D=\silverCodeFibres{c}{\lambda}$.
    Then $D$ is a function on $\lambda$ and
        for every $j\in\lambda$ we have
            $D(j)=\pow{\lambda}\times\funs{\lambda}{c(j)}$.
\end{proposition}
\begin{proof}
    $D$ is a relation by \auto{relation,silver_code_fibres}.
    $D$ is right-unique
        by \auto{rightunique,silver_code_fibres,pair_eq_iff}.
    $\dom{D}=\lambda$
        by \auto{setext,dom_iff,silver_code_fibres,pair_eq_iff}.
    Hence $D$ is a function on $\lambda$.
    For every $j\in\lambda$ we have
        $(j,\pow{\lambda}\times\funs{\lambda}{c(j)})\in D$
        by \obv{silver_code_fibres}.
    Hence for every $j\in\lambda$ we have
        $D(j)=\pow{\lambda}\times\funs{\lambda}{c(j)}$
        by \obv{function_apply_intro}.
\end{proof}

\begin{proposition}\label{silver_total_code_type}
    Suppose $S\subseteq\lambda$.
    Suppose $\emptyset\in B$.
    Suppose for every $u\in S$ we have $(E(u))(f(u))\in B$.
    Let $h=\silverTotalCode{E}{f}{S}{\lambda}$.
    Then $h\in\funs{\lambda}{B}$ and
        for every $v\in S$ we have $h(v)=(E(v))(f(v))$ and
        for every $w\in\lambda\setminus S$ we have $h(w)=\emptyset$.
\end{proposition}
\begin{proof}
    Let $p=\silverPartialCode{E}{f}{S}$.
    $p\in\funs{S}{B}$ by \auto{silver_partial_code_type}.
    $\identity{S}\in\funs{S}{S}$ by \obv{id_elem_funs}.
    Hence $\identity{S}\in\funs{S}{\lambda}$
        by \obv{funs_weaken_codom}.
    Show that $\identity{S}\in\inj{S}{\lambda}$.
    \begin{subproof}
        For all $x,y\in S$ such that
            $(\identity{S})(x)=(\identity{S})(y)$ we have $x=y$
            by \obv{id_apply}.
        Follows by \obv{inj}.
    \end{subproof}
    $h=\functionDomainExtension{\identity{S}}{p}{\lambda}{\emptyset}$
        by \obv{silver_total_code}.
    Hence $h\in\funs{\lambda}{B}$
        by \auto{function_domain_extension_type}.
    Show that for every $i\in S$ we have $h(i)=(E(i))(f(i))$.
    \begin{subproof}
        Fix $i\in S$.
        $(\identity{S})(i)=i$ by \obv{id_apply}.
        \begin{align*}
            h(i)
                &= (\functionDomainExtension{\identity{S}}{p}
                    {\lambda}{\emptyset})(i)
                \\
                &= (\functionDomainExtension{\identity{S}}{p}
                    {\lambda}{\emptyset})((\identity{S})(i))
                \\
                &= p(i)
                    \explanation{by
                        \obv{function_domain_extension_on_range}}
                \\
                &= (E(i))(f(i))
                    \explanation{by \auto{silver_partial_code_type}}
        \end{align*}
    \end{subproof}
    $\ran{\identity{S}}=S$ by \obv{id_ran}.
    For every $i\in\lambda\setminus S$ we have
        $h(i)=\emptyset$
        by \auto{function_domain_extension_off_range}.
\end{proof}

\begin{theorem}\label{silver_code_in_space}
    Let $\lambda$ be a regular cardinal.
    Suppose $\naturals\in\lambda$.
    Suppose $\normalMap{c}{\lambda}{\kappa}$.
    Suppose $c$ is index-dominating on $\lambda$.
    Suppose $T\subseteq\limitOrdinalsBelow{\lambda}$.
    Suppose $T$ is stationary in $\lambda$.
    Suppose $A$ is a function on $\lambda$.
    Suppose $E$ is a function on $T$.
    Suppose for every $u\in T$ we have
        $E(u)\in\inj{A(u)}{c(u)}$.
    Let $f\in\dependentProduct{\lambda}{A}$.
    Then $\silverCode{E}{c}{f}{T}{\lambda}
        \in\silverCodeSpace{c}{\lambda}$.
\end{theorem}
\begin{proof}
    Let $m=\silverRankIndex{E}{c}{f}{T}{\lambda}$.
    Let $S=\silverSupport{E}{c}{f}{T}{\lambda}$.
    $m\in\lambda$ and $S\subseteq T$ and
        $S$ is stationary in $\lambda$ and
        for every $i\in S$ we have $(E(i))(f(i))\in c(m)$
        by \obv{silver_support_properties}.
    $S\subseteq\lambda$ by \obv{stationary_in_elim}.
    $\emptyset\in c(m)$
        by \obv{index_dominating_normal_map_value_contains_zero}.
    Let $h=\silverTotalCode{E}{f}{S}{\lambda}$.
    $h\in\funs{\lambda}{c(m)}$ by \auto{silver_total_code_type}.
    $S\in\pow{\lambda}$ by \obv{powerset_intro}.
    Hence $(S,h)\in\pow{\lambda}\times\funs{\lambda}{c(m)}$
        by \obv{times_tuple_intro}.
    Let $D=\silverCodeFibres{c}{\lambda}$.
    $D$ is a function on $\lambda$ and
        $D(m)=\pow{\lambda}\times\funs{\lambda}{c(m)}$
        by \obv{silver_code_fibres_type}.
    Hence $(S,h)\in D(m)$.
    Hence $(m,(S,h))\in\indexedSum{\lambda}{D}$
        by \auto{indexed_sum_iff}.
    $\silverCode{E}{c}{f}{T}{\lambda}=(m,(S,h))$
        by \obv{silver_code}.
    Follows by \obv{silver_code_space}.
\end{proof}

Equal codes force agreement on a stationary set, hence equality within
the eventually different family.

\begin{theorem}\label{silver_equal_codes_imply_equal_functions}
    Let $\lambda$ be a regular cardinal.
    Suppose $\naturals\in\lambda$.
    Suppose $\normalMap{c}{\lambda}{\kappa}$.
    Suppose $c$ is index-dominating on $\lambda$.
    Suppose $T\subseteq\limitOrdinalsBelow{\lambda}$.
    Suppose $T$ is stationary in $\lambda$.
    Suppose $A$ is a function on $\lambda$.
    Suppose $E$ is a function on $T$.
    Suppose for every $u\in T$ we have
        $E(u)\in\inj{A(u)}{c(u)}$.
    Let $f\in\dependentProduct{\lambda}{A}$.
    Let $g\in\dependentProduct{\lambda}{A}$.
    Suppose $\eventuallyDifferent{f}{g}{\lambda}$.
    Suppose $\silverCode{E}{c}{f}{T}{\lambda}
        =\silverCode{E}{c}{g}{T}{\lambda}$.
    Then $f=g$.
\end{theorem}
\begin{proof}
    $\lambda$ is a limit ordinal
        by \obv{infinite_cardinal_is_limit_ordinal}.
    Suppose not.
    Let $m=\silverRankIndex{E}{c}{f}{T}{\lambda}$.
    Let $n=\silverRankIndex{E}{c}{g}{T}{\lambda}$.
    Let $S=\silverSupport{E}{c}{f}{T}{\lambda}$.
    Let $U=\silverSupport{E}{c}{g}{T}{\lambda}$.
    Let $h=\silverTotalCode{E}{f}{S}{\lambda}$.
    Let $q=\silverTotalCode{E}{g}{U}{\lambda}$.
    $m\in\lambda$ and $S\subseteq T$ and
        $S$ is stationary in $\lambda$ and
        for every $i\in S$ we have $(E(i))(f(i))\in c(m)$
        by \obv{silver_support_properties}.
    $n\in\lambda$ and $U\subseteq T$ and
        $U$ is stationary in $\lambda$ and
        for every $i\in U$ we have $(E(i))(g(i))\in c(n)$
        by \obv{silver_support_properties}.
    $S\subseteq\lambda$ and $U\subseteq\lambda$
        by \obv{stationary_in_elim}.
    $\emptyset\in c(m)$ and $\emptyset\in c(n)$
        by \obv{index_dominating_normal_map_value_contains_zero}.
    $h\in\funs{\lambda}{c(m)}$ and
        for every $i\in S$ we have $h(i)=(E(i))(f(i))$
        by \auto{silver_total_code_type}.
    $q\in\funs{\lambda}{c(n)}$ and
        for every $i\in U$ we have $q(i)=(E(i))(g(i))$
        by \auto{silver_total_code_type}.
    $\silverCode{E}{c}{f}{T}{\lambda}=(m,(S,h))$ and
        $\silverCode{E}{c}{g}{T}{\lambda}=(n,(U,q))$
        by \obv{silver_code}.
    Hence $m=n$ and $S=U$ and $h=q$
        by \auto{pair_eq_iff}.
    Show that for every $i\in S$ we have $f(i)=g(i)$.
    \begin{subproof}
        Fix $i\in S$.
        Hence $i\in U$.
        Hence $h(i)=(E(i))(f(i))$ and
            $q(i)=(E(i))(g(i))$.
        Hence $(E(i))(f(i))=(E(i))(g(i))$.
        $i\in T$ by \obv{subseteq}.
        Hence $i\in\limitOrdinalsBelow{\lambda}$ by \obv{subseteq}.
        Hence $i\in\lambda$ by \obv{limit_ordinals_below}.
        Hence $E(i)\in\inj{A(i)}{c(i)}$.
        $f(i)\in A(i)$ by \obv{dependent_product_coordinate}.
        $g(i)\in A(i)$ by \obv{dependent_product_coordinate}.
        Follows by \obv{inj}.
    \end{subproof}
    Hence $\nonstationaryIn{S}{\lambda}$
        by \obv{eventually_different_agreement_set_nonstationary}.
    Contradiction by \auto{nonstationary_iff_not_stationary}.
\end{proof}

\begin{theorem}[The coding map is injective]\label{silver_code_map_is_injection}
    Let $\lambda$ be a regular cardinal.
    Suppose $\naturals\in\lambda$.
    Suppose $\normalMap{c}{\lambda}{\kappa}$.
    Suppose $c$ is index-dominating on $\lambda$.
    Suppose $T\subseteq\limitOrdinalsBelow{\lambda}$.
    Suppose $T$ is stationary in $\lambda$.
    Suppose $A$ is a function on $\lambda$.
    Suppose $E$ is a function on $T$.
    Suppose for every $u\in T$ we have
        $E(u)\in\inj{A(u)}{c(u)}$.
    Suppose $F\subseteq\dependentProduct{\lambda}{A}$.
    Suppose $\eventuallyDifferentFamily{F}{\lambda}$.
    Then $\silverCodeMap{E}{c}{F}{T}{\lambda}
        \in\inj{F}{\silverCodeSpace{c}{\lambda}}$.
\end{theorem}
\begin{proof}
    Let $H=\silverCodeMap{E}{c}{F}{T}{\lambda}$.
    $H$ is a relation by \auto{relation,silver_code_map}.
    $H$ is right-unique
        by \auto{rightunique,silver_code_map,pair_eq_iff}.
    $\dom{H}=F$
        by \auto{setext,dom_iff,silver_code_map,pair_eq_iff}.
    Show that for every $f\in F$ we have
        $H(f)\in\silverCodeSpace{c}{\lambda}$.
    \begin{subproof}
        Fix $f\in F$.
        Hence $f\in\dependentProduct{\lambda}{A}$ by \obv{subseteq}.
        $\silverCode{E}{c}{f}{T}{\lambda}
            \in\silverCodeSpace{c}{\lambda}$
            by \obv{silver_code_in_space}.
        $(f,\silverCode{E}{c}{f}{T}{\lambda})\in H$
            by \obv{silver_code_map}.
        Hence $H(f)=\silverCode{E}{c}{f}{T}{\lambda}$
            by \obv{function_apply_intro}.
    \end{subproof}
    Hence $H$ is a function to $\silverCodeSpace{c}{\lambda}$.
    Hence $H\in\funs{F}{\silverCodeSpace{c}{\lambda}}$
        by \obv{funs_intro}.
    Show that $H$ is injective.
    \begin{subproof}
        Show that for all $f,g\in\dom{H}$ such that $H(f)=H(g)$ we have $f=g$.
        \begin{subproof}
            Fix $f,g\in\dom{H}$ such that $H(f)=H(g)$.
            Hence $f,g\in F$.
            Hence $f,g\in\dependentProduct{\lambda}{A}$
                by \obv{subseteq}.
            $H(f)=\silverCode{E}{c}{f}{T}{\lambda}$ and
                $H(g)=\silverCode{E}{c}{g}{T}{\lambda}$
                by \auto{function_apply_intro,silver_code_map}.
            Hence $\silverCode{E}{c}{f}{T}{\lambda}
                =\silverCode{E}{c}{g}{T}{\lambda}$.
            \begin{byCase}
                \caseOf{$f=g$.}
                    Follows.
                \caseOf{$f\neq g$.}
                    Hence $\eventuallyDifferent{f}{g}{\lambda}$
                        by \obv{eventually_different_family}.
                    Follows by
                        \obv{silver_equal_codes_imply_equal_functions}.
            \end{byCase}
        \end{subproof}
        Follows by \auto{injective_function}.
    \end{subproof}
    Hence $H\in\inj{F}{\silverCodeSpace{c}{\lambda}}$
        by \auto{inj,injective_function}.
\end{proof}

\subsection{Counting the codes}

\begin{theorem}\label{silver_code_fibre_cardinal_bound}
    Let $\kappa$ be an infinite cardinal.
    Suppose $\ordinalMap{c}{\lambda}{\kappa}$.
    Suppose $\cardinalSequenceValues{c}{\lambda}$.
    Suppose $\powerClosedBelow{\lambda}{\kappa}$.
    Let $j\in\lambda$.
    Then $\cardLe{\pow{\lambda}\times\funs{\lambda}{c(j)}}{\kappa}$.
\end{theorem}
\begin{proof}
    $\cardOne\in\naturals$
        by \auto{cardinal_one,omega_empty,omega_successor}.
    Hence $\cardTwo\in\naturals$
        by \auto{cardinal_two_successor,omega_successor}.
    $\naturals\subseteq\kappa$
        by \obv{naturals_subseteq_infinite_cardinal}.
    Hence $\cardTwo\in\kappa$ by \obv{subseteq}.
    Hence $\card{\pow{\lambda}}\in\kappa$
        by \obv{power_closed_below_powerset}.
    Hence $\cardLe{\pow{\lambda}}{\kappa}$
        by \obv{card_belongs_to_cardinal_implies_cardinal_le}.
    $c\in\funs{\lambda}{\kappa}$ by \obv{ordinal_map}.
    Hence $c(j)\in\kappa$ by \obv{funs_type_apply}.
    Hence $\cardPower{c(j)}{\lambda}\in\kappa$
        by \obv{power_closed_below}.
    $\card{\funs{\lambda}{c(j)}}=\cardPower{c(j)}{\lambda}$
        by \obv{cardinal_power}.
    Hence $\card{\funs{\lambda}{c(j)}}\in\kappa$.
    Hence $\cardLe{\funs{\lambda}{c(j)}}{\kappa}$
        by \obv{card_belongs_to_cardinal_implies_cardinal_le}.
    Take $e$ such that $e\in\inj{\pow{\lambda}}{\kappa}$
        by \auto{cardinal_le}.
    Take $d$ such that $d\in\inj{\funs{\lambda}{c(j)}}{\kappa}$
        by \auto{cardinal_le}.
    Take $q$ such that
        $q\in\inj{\pow{\lambda}\times\funs{\lambda}{c(j)}}
            {\kappa\times\kappa}$
        by \auto{times_injection_congruent}.
    Hence $\cardLe{\pow{\lambda}\times\funs{\lambda}{c(j)}}
        {\kappa\times\kappa}$ by \auto{cardinal_le}.
    $\kappa\times\kappa\approx\kappa$
        by \obv{infinite_cardinal_square}.
    Hence $\cardLe{\kappa\times\kappa}{\kappa}$
        by \auto{equinum_cardinal_le}.
    Follows by \obv{cardinal_le_transitive}.
\end{proof}

\begin{theorem}[Counting the codes]\label{silver_code_space_cardinal_bound}
    Let $\kappa$ be an infinite cardinal.
    Suppose $\lambda\in\kappa$.
    Suppose $\ordinalMap{c}{\lambda}{\kappa}$.
    Suppose $\cardinalSequenceValues{c}{\lambda}$.
    Suppose $\powerClosedBelow{\lambda}{\kappa}$.
    Then $\cardLe{\silverCodeSpace{c}{\lambda}}{\kappa}$.
\end{theorem}
\begin{proof}
    $\lambda\subseteq\kappa$
        by \obv{ordinal_elem_implies_subseteq}.
    Hence $\cardLe{\lambda}{\kappa}$
        by \obv{subseteq_cardinal_le}.
    Let $D=\silverCodeFibres{c}{\lambda}$.
    $D$ is a function on $\lambda$ and
        for every $j\in\lambda$ we have
            $D(j)=\pow{\lambda}\times\funs{\lambda}{c(j)}$
        by \obv{silver_code_fibres_type}.
    For every $j\in\lambda$ we have $\cardLe{D(j)}{\kappa}$
        by \auto{silver_code_fibre_cardinal_bound}.
    Hence $\cardLe{\indexedSum{\lambda}{D}}{\kappa}$
        by \obv{indexed_sum_infinite_cardinal_bound}.
    Follows by \auto{silver_code_space}.
\end{proof}

\subsection{The first bound}

\begin{theorem}[The first bound for eventually different families]\label{silver_first_almost_disjoint_bound}
    Let $\kappa$ be a singular cardinal.
    Let $\lambda=\cofinality{\kappa}$.
    Suppose $\naturals\in\lambda$.
    Suppose $\normalMap{c}{\lambda}{\kappa}$.
    Suppose $\cardinalSequenceValues{c}{\lambda}$.
    Suppose $c$ is index-dominating on $\lambda$.
    Suppose $\powerClosedBelow{\lambda}{\kappa}$.
    Suppose $T$ is stationary in $\lambda$.
    Suppose $A$ is a function on $\lambda$.
    Suppose $F\subseteq\dependentProduct{\lambda}{A}$.
    Suppose $\eventuallyDifferentFamily{F}{\lambda}$.
    Suppose for every $u\in T$ we have $\cardLe{A(u)}{c(u)}$.
    Then $\cardLe{F}{\kappa}$.
\end{theorem}
\begin{proof}
    $\kappa$ is a limit ordinal
        by \obv{infinite_cardinal_is_limit_ordinal}.
    $\cofinality{\kappa}$ is a regular cardinal
        by \obv{cofinality_of_limit_is_regular}.
    Hence $\lambda$ is a regular cardinal.
    $\ordinalMap{c}{\lambda}{\kappa}$ by \obv{normal_map}.
    $\cofinality{\kappa}\in\kappa$
        by \obv{cofinality_belongs_to_singular_cardinal}.
    Hence $\lambda\in\kappa$.
    $\silverLimitSupport{T}{\lambda}$ is stationary in $\lambda$
        by \obv{silver_limit_support_is_stationary}.
    Let $U=\silverLimitSupport{T}{\lambda}$.
    Hence $U$ is stationary in $\lambda$.
    Hence $U\subseteq\lambda$ by \obv{stationary_in_elim}.
    $U\subseteq T$ and $U\subseteq\limitOrdinalsBelow{\lambda}$
        by \auto{silver_limit_support,subseteq,inter}.
    Let $A_0=\restrl{A}{U}$.
    Let $c_0=\restrl{c}{U}$.
    $A$ is a function and $\dom{A}=\lambda$.
    Hence $\dom{A_0}=U$
        by \auto{dom_of_restrl_eq_inter,inter_absorb_supseteq_left}.
    Hence $A_0$ is a function on $U$.
    $c\in\funs{\lambda}{\kappa}$
        by \obv{normal_map_type}.
    Hence $c$ is a function and $\dom{c}=\lambda$
        by \auto{funs_is_function,funs}.
    Hence $\dom{c_0}=U$
        by \auto{dom_of_restrl_eq_inter,inter_absorb_supseteq_left}.
    Hence $c_0$ is a function on $U$.
    For every $i\in U$ we have $A_0(i)=A(i)$ and $c_0(i)=c(i)$
        by \auto{restrl_of_function_apply}.
    For every $i\in U$ we have $\cardLe{A_0(i)}{c_0(i)}$
        by \auto{subseteq}.
    Take $E$ such that $E$ is a function on $U$ and
        for every $i\in U$ we have
            $E(i)\in\inj{A_0(i)}{c_0(i)}$
        by \obv{indexed_family_variable_injections}.
    Hence for every $i\in U$ we have
        $E(i)\in\inj{A(i)}{c(i)}$ by \auto.
    $\silverCodeMap{E}{c}{F}{U}{\lambda}
        \in\inj{F}{\silverCodeSpace{c}{\lambda}}$
        by \obv{silver_code_map_is_injection}.
    Hence $\cardLe{F}{\silverCodeSpace{c}{\lambda}}$
        by \auto{cardinal_le}.
    $\cardLe{\silverCodeSpace{c}{\lambda}}{\kappa}$
        by \obv{silver_code_space_cardinal_bound}.
    Follows by \obv{cardinal_le_transitive}.
\end{proof}
